% Generated by the blueprint tool (blueprint-pdf/1) from a blueprint snapshot; edit the
% Markdown sources, not this file.  Compile with `latexmk -pdf` or by running
% `pdflatex` twice (the second run settles the contents and the page references).
% Images, if any, are read from assets/ next to this file.
\documentclass[11pt]{article}
\usepackage[T1]{fontenc}
\usepackage[utf8]{inputenc}
\usepackage{lmodern}
\usepackage{textcomp}
\usepackage[a4paper,margin=2.5cm,headheight=14pt]{geometry}
\usepackage{amsmath,amssymb,amsthm}
\usepackage{amscd}  % KaTeX's CD diagrams; part of amsmath, so always there
\usepackage{mathtools}
\usepackage{bm}
\usepackage{graphicx}
\usepackage{xcolor}
\usepackage{array,booktabs,tabularx}
\usepackage{fancyvrb}
\usepackage[normalem]{ulem}
\usepackage{microtype}
% Optional, and only with their fonts: a small TeX installation has the
% packages without the fonts, which fails at the first formula.
\IfFileExists{mathrsfs.sty}{\IfFileExists{rsfs10.tfm}{\usepackage{mathrsfs}}{}}{}
\IfFileExists{stmaryrd.sty}{\IfFileExists{stmary10.tfm}{\usepackage{stmaryrd}}{}}{}
\IfFileExists{cancel.sty}{\usepackage{cancel}}{}
\usepackage[hyphens]{url}
\usepackage[pdfusetitle,bookmarksnumbered=false]{hyperref}
\usepackage{bookmark}
\definecolor{bplink}{HTML}{1F4E9C}
\definecolor{bpmuted}{HTML}{5F676E}
\hypersetup{colorlinks=true,linkcolor=bplink,urlcolor=bplink,citecolor=bplink}
\setlength{\parindent}{0pt}
\setlength{\parskip}{0.55em plus 0.2em minus 0.1em}
% Lists as tight as the page's.
\makeatletter
\def\@listi{\leftmargin\leftmargini\parsep 0pt\topsep 0.3em plus 0.1em\itemsep 0.1em plus 0.1em}
\let\@listI\@listi
\def\@listii{\leftmargin\leftmarginii\labelwidth\leftmarginii\advance\labelwidth-\labelsep\parsep 0pt\topsep 0.1em\itemsep 0.1em}
% An ordered list that does not start at 1.
\newcommand{\bpstartat}[1]{\setcounter{\@enumctr}{\numexpr#1-1\relax}}
\makeatother
\setcounter{tocdepth}{3}
\emergencystretch=3em
\allowdisplaybreaks

% A numbered statement, set as a paper sets one: "Theorem 2.1 (Title)."
% #1 is the lead ("Theorem~2.1"), #2 what follows it in roman.
\newenvironment{bpstatement}[2]{\par\addvspace{0.6\baselineskip}\noindent\textbf{#1}#2\textbf{.}\enspace\ignorespaces}
  {\par\addvspace{0.4\baselineskip}}
\newenvironment{bpstatementit}[2]{\begin{bpstatement}{#1}{#2}\itshape\ignorespaces}{\end{bpstatement}}
% A step: a relation with prose of its own, read under the object it starts from.
\newenvironment{bpstep}[1]{\par\addvspace{0.3\baselineskip}\begin{list}{}{\leftmargin=1.6em\rightmargin=0pt\topsep=0pt}\item\relax
  {\small\color{bpmuted}#1}\par\nopagebreak}{\end{list}\par}
% Notes about where an object is written out.
\newenvironment{bpnote}{\par\small\color{bpmuted}\itshape}{\par}
\newcommand{\bpomitted}[1]{\par{\small\color{bpmuted}\itshape #1\par}}
\newcommand{\bpbodyheading}[1]{\par\addvspace{0.4\baselineskip}\noindent\textbf{#1}\par\nopagebreak}
\newcommand{\bprule}{\par\medskip\noindent\hfil\rule{0.3\linewidth}{0.4pt}\hfil\par\medskip}
\newcommand{\bpboundary}[1]{\par\nopagebreak{\small\color{bpmuted}#1\par}\nopagebreak}
% A reference to something this export leaves out, and one that resolves nowhere.
\newcommand{\bpout}[1]{#1\textsuperscript{\hyperlink{bp-outside}{\ensuremath{\ast}}}}
\newcommand{\bpbroken}[1]{\textcolor{red!70!black}{[#1?]}}
\newcommand{\bpslug}[1]{\texttt{#1}}
% A bibliography entry: label, anchor, entry; the label hangs in the margin.
\newcommand{\bpbibitem}[3]{\par\setbox0\hbox{[#1]\enspace}\ifdim\wd0<5em\setbox0\hbox to 5em{[#1]\hfil}\fi
  \hangindent=\wd0 \hangafter=1 \noindent\box0\phantomsection\label{#2}#3\par\smallskip}
% A character the export knows no way to set.
\newcommand{\bpglyph}[1]{\fbox{\scriptsize #1}}
% Images, at their own size unless they are wider than the line.
\newsavebox{\bpimagebox}
\newcommand{\bpimage}[1]{\sbox{\bpimagebox}{\includegraphics{#1}}%
  \ifdim\wd\bpimagebox>\linewidth\resizebox{\linewidth}{!}{\usebox{\bpimagebox}}\else\usebox{\bpimagebox}\fi}
\newcommand{\bpmissingimage}[1]{\fbox{\parbox{0.8\linewidth}{\centering\small\color{bpmuted}Image not included: #1}}}
\DefineVerbatimEnvironment{bpverbatim}{Verbatim}{fontsize=\small,xleftmargin=1em}
\renewcommand{\qedsymbol}{\ensuremath{\blacksquare}}
\hypersetup{pdftitle={Weil Conjectures — Appendices},pdfsubject={Part Appendices; all levels},pdfcreator={blueprint}}
% KaTeX commands the website knows
\providecommand{\R}{\mathbb{R}}\providecommand{\N}{\mathbb{N}}\providecommand{\Z}{\mathbb{Z}}
\providecommand{\Q}{\mathbb{Q}}\providecommand{\C}{\mathbb{C}}
\providecommand{\Reals}{\mathbb{R}}\providecommand{\reals}{\mathbb{R}}\providecommand{\Complex}{\mathbb{C}}
\providecommand{\cnums}{\mathbb{C}}\providecommand{\natnums}{\mathbb{N}}\providecommand{\Bbb}[1]{\mathbb{#1}}
\providecommand{\bold}[1]{\mathbf{#1}}\providecommand{\frak}[1]{\mathfrak{#1}}\providecommand{\mathbbm}[1]{\mathbb{#1}}
\providecommand{\lang}{\langle}\providecommand{\rang}{\rangle}\providecommand{\bra}[1]{\langle #1\rvert}
\providecommand{\ket}[1]{\lvert #1\rangle}\providecommand{\braket}[1]{\langle #1\rangle}
\providecommand{\set}[1]{\{#1\}}\providecommand{\Set}[1]{\left\{#1\right\}}
\providecommand{\infin}{\infty}\providecommand{\isin}{\in}\providecommand{\plusmn}{\pm}\providecommand{\sub}{\subset}
\providecommand{\sube}{\subseteq}\providecommand{\supe}{\supseteq}\providecommand{\exist}{\exists}
\providecommand{\bull}{\bullet}\providecommand{\larr}{\leftarrow}\providecommand{\rarr}{\rightarrow}
\providecommand{\lrarr}{\leftrightarrow}\providecommand{\Larr}{\Leftarrow}\providecommand{\Rarr}{\Rightarrow}
\providecommand{\Lrarr}{\Leftrightarrow}\providecommand{\uarr}{\uparrow}\providecommand{\darr}{\downarrow}
\providecommand{\harr}{\leftrightarrow}\providecommand{\lArr}{\Leftarrow}\providecommand{\rArr}{\Rightarrow}
\providecommand{\hArr}{\Leftrightarrow}\providecommand{\alef}{\aleph}\providecommand{\alefsym}{\aleph}
\providecommand{\real}{\Re}\providecommand{\image}{\Im}\providecommand{\weierp}{\wp}
\providecommand{\Alpha}{\mathrm{A}}\providecommand{\Beta}{\mathrm{B}}\providecommand{\Epsilon}{\mathrm{E}}
\providecommand{\Zeta}{\mathrm{Z}}\providecommand{\Eta}{\mathrm{H}}\providecommand{\Iota}{\mathrm{I}}
\providecommand{\Kappa}{\mathrm{K}}\providecommand{\Mu}{\mathrm{M}}\providecommand{\Nu}{\mathrm{N}}
\providecommand{\Omicron}{\mathrm{O}}\providecommand{\Rho}{\mathrm{P}}\providecommand{\Tau}{\mathrm{T}}
\providecommand{\Chi}{\mathrm{X}}\providecommand{\omicron}{o}\providecommand{\gt}{>}\providecommand{\lt}{<}
\providecommand{\And}{\&}\providecommand{\hdashline}{\hline}\providecommand{\operatornamewithlimits}[1]{\operatorname*{#1}}
\providecommand{\llbracket}{[\![}\providecommand{\rrbracket}{]\!]}\providecommand{\cancel}[1]{#1}
\providecommand{\mathscr}{\mathcal}
% the characters beyond ASCII this document uses
\DeclareUnicodeCharacter{2160}{\textup{I}}
\DeclareUnicodeCharacter{2161}{\textup{II}}
\DeclareUnicodeCharacter{2162}{\textup{III}}
% the project's macros ([katex.macros] in blueprint.toml)
\providecommand{\A}{}\renewcommand{\A}{\mathbf A}
\providecommand{\AEn}{}\renewcommand{\AEn}{A_{E,n}}
\providecommand{\Aaf}{}\renewcommand{\Aaf}{\mathbf A}
\providecommand{\Ab}{}\renewcommand{\Ab}{\mathsf{Ab}}
\providecommand{\AffSch}{}\renewcommand{\AffSch}{\operatorname{AffSch}}
\providecommand{\Aone}{}\renewcommand{\Aone}{\mathbf A^1}
\providecommand{\Bl}{}\renewcommand{\Bl}{\operatorname{Bl}}
\providecommand{\Br}{}\renewcommand{\Br}{\mathrm{Br}}
\providecommand{\C}{}\renewcommand{\C}{\mathbb C}
\providecommand{\Cech}{}\renewcommand{\Cech}{\check{H}}
\providecommand{\Coext}{}\renewcommand{\Coext}{\operatorname{Coext}}
\providecommand{\Cone}{}\renewcommand{\Cone}{\operatorname{Cone}}
\providecommand{\Cons}{}\renewcommand{\Cons}{\operatorname{Cons}}
\providecommand{\Cref}{}\renewcommand{\Cref}[1]{\text{#1}}
\providecommand{\D}{}\renewcommand{\D}{D}
\providecommand{\DLAFT}{}\renewcommand{\DLAFT}{\operatorname{FT}}
\providecommand{\DLAFr}{}\renewcommand{\DLAFr}{\operatorname{Fr}}
\providecommand{\DLAQ}{}\renewcommand{\DLAQ}{{\overline{\mathbf Q}_\ell}}
\providecommand{\DLRQ}{}\renewcommand{\DLRQ}{{\overline{\mathbf Q}_\ell}}
\providecommand{\Dbc}{}\renewcommand{\Dbc}{D^b_c}
\providecommand{\Dc}{}\renewcommand{\Dc}{D_{\mathrm{cons}}}
\providecommand{\Dctf}{}\renewcommand{\Dctf}{D^b_{ctf}}
\providecommand{\Dic}{}\renewcommand{\Dic}{D_{\IndCons}}
\providecommand{\Dplus}{}\renewcommand{\Dplus}{D^{+}}
\providecommand{\End}{}\renewcommand{\End}{\operatorname{End}}
\providecommand{\Et}{}\renewcommand{\Et}{\mathrm{\acute{e}t}}
\providecommand{\Ex}{}\renewcommand{\Ex}{\mathrm{Ex}}
\providecommand{\Ext}{}\renewcommand{\Ext}{\operatorname{Ext}}
\providecommand{\F}{}\renewcommand{\F}{\mathbb F}
\providecommand{\FT}{}\renewcommand{\FT}{\operatorname{FT}}
\providecommand{\Fbar}{}\renewcommand{\Fbar}{{\overline{\mathbf F}_q}}
\providecommand{\Fell}{}\renewcommand{\Fell}{\mathbf F_\ell}
\providecommand{\Fib}{}\renewcommand{\Fib}{\operatorname{Fib}}
\providecommand{\Fix}{}\renewcommand{\Fix}{\operatorname{Fix}}
\providecommand{\Fp}{}\renewcommand{\Fp}{{\mathbf F_{p}}}
\providecommand{\Fq}{}\renewcommand{\Fq}{{\mathbf F_q}}
\providecommand{\Fqn}{}\renewcommand{\Fqn}{{\mathbf F_{q^{n}}}}
\providecommand{\Frob}{}\renewcommand{\Frob}{\operatorname{Frob}}
\providecommand{\GL}{}\renewcommand{\GL}{\operatorname{GL}}
\providecommand{\Gal}{}\renewcommand{\Gal}{\operatorname{Gal}}
\providecommand{\Hens}{}\renewcommand{\Hens}{\operatorname{Hens}}
\providecommand{\Hilb}{}\renewcommand{\Hilb}{\operatorname{Hilb}}
\providecommand{\Hom}{}\renewcommand{\Hom}{\operatorname{Hom}}
\providecommand{\IndCons}{}\renewcommand{\IndCons}{\operatorname{IndCons}}
\providecommand{\Kar}{}\renewcommand{\Kar}{\operatorname{Kar}}
\providecommand{\LT}{}\renewcommand{\LT}{\operatorname{LT}}
\providecommand{\Lie}{}\renewcommand{\Lie}{\operatorname{Lie}}
\providecommand{\Loc}{}\renewcommand{\Loc}{\operatorname{Loc}}
\providecommand{\Mod}{}\renewcommand{\Mod}{\mathsf{Mod}}
\providecommand{\N}{}\renewcommand{\N}{\mathbb{N}}
\providecommand{\Needspace}{}\renewcommand{\Needspace}{\par}
\providecommand{\O}{}\renewcommand{\O}{{\mathcal{O}}}
\providecommand{\Ob}{}\renewcommand{\Ob}{\mathrm{Ob}}
\providecommand{\Oe}{}\renewcommand{\Oe}{\mathcal O_E}
\providecommand{\PP}{}\renewcommand{\PP}{\mathbf{P}}
\providecommand{\Perf}{}\renewcommand{\Perf}{\operatorname{Perf}}
\providecommand{\Pic}{}\renewcommand{\Pic}{\mathrm{Pic}}
\providecommand{\Pm}{}\renewcommand{\Pm}{\mathbf P^{m}}
\providecommand{\Pone}{}\renewcommand{\Pone}{\mathbf P^{1}}
\providecommand{\Pqn}{}\renewcommand{\Pqn}{\mathbf P^{n}}
\providecommand{\ProFin}{}\renewcommand{\ProFin}{\operatorname{ProFin}}
\providecommand{\Proet}{}\renewcommand{\Proet}{\mathrm{pro\acute{e}t}}
\providecommand{\Q}{}\renewcommand{\Q}{\mathbb Q}
\providecommand{\Qb}{}\renewcommand{\Qb}{\overline{\mathbf Q}_{\ell}}
\providecommand{\Qbar}{}\renewcommand{\Qbar}{\overline{\mathbf Q}}
\providecommand{\QbarX}{}\renewcommand{\QbarX}{\overline{\mathbf Q}_{\ell,X}}
\providecommand{\Qbarl}{}\renewcommand{\Qbarl}{{\overline{\mathbf Q}_{\ell}}}
\providecommand{\Qell}{}\renewcommand{\Qell}{{\mathbf{Q}_{\ell}}}
\providecommand{\Ql}{}\renewcommand{\Ql}{\mathbf Q_{\ell}}
\providecommand{\Qlb}{}\renewcommand{\Qlb}{{\overline{\mathbf Q}_\ell}}
\providecommand{\Qlbar}{}\renewcommand{\Qlbar}{\overline{\mathbf Q}_\ell}
\providecommand{\R}{}\renewcommand{\R}{\mathbb{R}}
\providecommand{\RG}{}\renewcommand{\RG}{\operatorname{\mathrm{R}\Gamma}}
\providecommand{\RGamma}{}\renewcommand{\RGamma}{R\Gamma}
\providecommand{\RGammac}{}\renewcommand{\RGammac}{R\Gamma_c}
\providecommand{\RHom}{}\renewcommand{\RHom}{R\mathcal{H}om}
\providecommand{\RHomX}{}\renewcommand{\RHomX}{R\mathcal H om_X}
\providecommand{\Rat}{}\renewcommand{\Rat}{\operatorname{IsRat}}
\providecommand{\Res}{}\renewcommand{\Res}{\operatorname{Res}}
\providecommand{\Rlim}{}\renewcommand{\Rlim}{R\!\varprojlim}
\providecommand{\SIm}{}\renewcommand{\SIm}{\operatorname{SIm}}
\providecommand{\Sch}{}\renewcommand{\Sch}{\mathsf{Sch}}
\providecommand{\Sh}{}\renewcommand{\Sh}{\operatorname{Sh}}
\providecommand{\Span}{}\renewcommand{\Span}{\operatorname{span}}
\providecommand{\Spec}{}\renewcommand{\Spec}{\operatorname{Spec}}
\providecommand{\Tr}{}\renewcommand{\Tr}{\operatorname{Tr}}
\providecommand{\Trsm}{}\renewcommand{\Trsm}{\operatorname{Tr}^{\mathrm{sm}}}
\providecommand{\Wt}{}\renewcommand{\Wt}{\operatorname{Wt}}
\providecommand{\Z}{}\renewcommand{\Z}{\mathbb{Z}}
\providecommand{\Zar}{}\renewcommand{\Zar}{\mathrm{zar}}
\providecommand{\Zell}{}\renewcommand{\Zell}{{\mathbf{Z}_{\ell}}}
\providecommand{\Zl}{}\renewcommand{\Zl}{{\mathbf Z_\ell}}
\providecommand{\affet}{}\renewcommand{\affet}[1]{{#1}^{\mathrm{aff}}_{\mathrm{ét}}}
\providecommand{\affproet}{}\renewcommand{\affproet}[1]{{#1}^{\mathrm{aff}}_{\mathrm{proét}}}
\providecommand{\boxt}{}\renewcommand{\boxt}{\boxtimes_A^{\mathrm L}}
\providecommand{\cF}{}\renewcommand{\cF}{\mathcal{F}}
\providecommand{\cHom}{}\renewcommand{\cHom}{R\mathcal{H}\!om}
\providecommand{\cI}{}\renewcommand{\cI}{\mathcal{I}}
\providecommand{\cL}{}\renewcommand{\cL}{\mathcal L}
\providecommand{\cM}{}\renewcommand{\cM}{\mathcal M}
\providecommand{\cN}{}\renewcommand{\cN}{\mathcal N}
\providecommand{\cO}{}\renewcommand{\cO}{\mathcal O}
\providecommand{\cT}{}\renewcommand{\cT}{\mathcal T}
\providecommand{\checked}{}\renewcommand{\checked}{\color{red}\aftergroup\restorecolor}
\providecommand{\coeff}{}\renewcommand{\coeff}{\operatorname{coeff}}
\providecommand{\coker}{}\renewcommand{\coker}{\operatorname{coker}}
\providecommand{\colim}{}\renewcommand{\colim}{\operatorname{colim}}
\providecommand{\constructedby}{}\renewcommand{\constructedby}[1]{\par\smallskip\noindent\textcolor{purple!60!black}{\textsf{Constructed by:} }
  \footnotesize\sloppy\nolinkurl{#1}\par}
\providecommand{\deepresult}{}\renewcommand{\deepresult}[1]{\par\smallskip\noindent\textcolor{orange!45!black}{
  \textsf{Blueprint role: named deep theorem node; see the cluster proof architecture.}}
  \par\noindent\textcolor{orange!45!black}{\textsf{Mathematical source:} }#1\par}
\providecommand{\dt}{}\renewcommand{\dt}{\mathbin{\otimes^{\mathbf L}}}
\providecommand{\dw}{}\renewcommand{\dw}{\operatorname{dw}}
\providecommand{\et}{}\renewcommand{\et}[1]{{#1}_{\mathrm{ét}}}
\providecommand{\ev}{}\renewcommand{\ev}{\operatorname{ev}}
\providecommand{\fib}{}\renewcommand{\fib}{\operatorname{fib}}
\providecommand{\formalizationnote}{}\renewcommand{\formalizationnote}[1]{\begin{quote}\small\raggedright\textbf{Formalization note.} #1\end{quote}}
\providecommand{\fpqc}{}\renewcommand{\fpqc}{X_{\mathrm{fpqc}}}
\providecommand{\gammacont}{}\renewcommand{\gammacont}{\Gamma_{\mathrm{cont}}}
\providecommand{\gm}{}\renewcommand{\gm}{\mathbf G_m}
\providecommand{\hocolim}{}\renewcommand{\hocolim}{\operatorname{hocolim}}
\providecommand{\holim}{}\renewcommand{\holim}{\operatorname{holim}}
\providecommand{\hrB}{}\renewcommand{\hrB}{\@ifnextchar]{\hrEnd}{\hrStep}}
\providecommand{\hrEnd}{}\renewcommand{\hrEnd}{\@firstoftwo{}}
\providecommand{\hrStep}{}\renewcommand{\hrStep}{\@firstoftwo{\hrB}}
\providecommand{\id}{}\renewcommand{\id}{\operatorname{id}}
\providecommand{\kunFl}{}\renewcommand{\kunFl}{\mathbf F_\ell}
\providecommand{\kunZn}{}\renewcommand{\kunZn}{\mathbf Z/n\mathbf Z}
\providecommand{\kuncd}{}\renewcommand{\kuncd}{\operatorname{cd}}
\providecommand{\kundim}{}\renewcommand{\kundim}{\dim}
\providecommand{\kunotimes}{}\renewcommand{\kunotimes}[1]{\otimes^{\mathbf L}_{#1}}
\providecommand{\kunsupp}{}\renewcommand{\kunsupp}{\operatorname{Supp}}
\providecommand{\lotimes}{}\renewcommand{\lotimes}{\otimes^{\mathrm L}}
\providecommand{\m}{}\renewcommand{\m}{{\mathfrak{m}}}
\providecommand{\n}{}\renewcommand{\n}{{\mathfrak{n}}}
\providecommand{\nagataGap}{}\renewcommand{\nagataGap}[1]{\par\noindent\textbf{Formalization gap.} #1\par}
\providecommand{\openobligation}{}\renewcommand{\openobligation}[1]{\par\smallskip\noindent\textcolor{red!65!black}{\textsf{Open obligation:} }
  \footnotesize\sloppy\nolinkurl{#1}\par}
\providecommand{\p}{}\renewcommand{\p}{{\mathfrak{p}}}
\providecommand{\pr}{}\renewcommand{\pr}{\operatorname{pr}}
\providecommand{\preshv}{}\renewcommand{\preshv}[1]{\mathrm{PreShv}(#1)}
\providecommand{\pro}{}\renewcommand{\pro}[1]{\mathrm{Pro}(#1)}
\providecommand{\proet}{}\renewcommand{\proet}{X_{\mathrm{proét}}}
\providecommand{\q}{}\renewcommand{\q}{{\mathfrak{q}}}
\providecommand{\rad}{}\renewcommand{\rad}{\operatorname{rad}}
\providecommand{\rank}{}\renewcommand{\rank}{\operatorname{rank}}
\providecommand{\red}{}\renewcommand{\red}{\color{red}}
\providecommand{\restorecolor}{}\renewcommand{\restorecolor}{\color{black}}
\providecommand{\retiredobligation}{}\renewcommand{\retiredobligation}[1]{\par\smallskip\noindent\textcolor{red!75!black}{\textsf{Retired invalid route:} }
  \footnotesize\sloppy\nolinkurl{#1}\par}
\providecommand{\rk}{}\renewcommand{\rk}{\operatorname{rk}}
\providecommand{\sep}{}\renewcommand{\sep}{\textnormal{sep}}
\providecommand{\shv}{}\renewcommand{\shv}[1]{\mathrm{Shv}(#1)}
\providecommand{\spec}{}\renewcommand{\spec}[1]{\mathrm{Spec}(#1)}
\providecommand{\stacks}{}\renewcommand{\stacks}[1]{\href{https://stacks.math.columbia.edu/tag/#1}{\textnormal{Stacks, Tag~#1}}}
\providecommand{\stacksproject}{}\renewcommand{\stacksproject}[1]{{\cite[\href{https://stacks.math.columbia.edu/tag/#1}{Tag #1}]{stacks-project}}}
\providecommand{\statusnote}{}\renewcommand{\statusnote}[1]{\begin{quote}\small\raggedright\textbf{Status.} #1\end{quote}}
\providecommand{\tildeideal}{}\renewcommand{\tildeideal}[2]{\widetilde{(#1, #2)}}
\providecommand{\tildering}{}\renewcommand{\tildering}[3]{{#1}_{\widetilde{(#2, #3)}}}
\providecommand{\tr}{}\renewcommand{\tr}{\operatorname{tr}}
\providecommand{\wt}{}\renewcommand{\wt}{\operatorname{wt}}

\begin{document}
\pagestyle{plain}
\begin{center}
{\LARGE\bfseries Weil Conjectures\par}
\vspace{0.7em}
{\large Part: Appendices\par}
\vspace{0.5em}
{\small\color{bpmuted}Appendices left out\par}
{\small\color{bpmuted}Snapshot 2026-09-26, revision 8e4cb88e2a\par}
\end{center}
\bigskip

% object: part-Appendices (blueprint/part-Appendices/_section.md)
\part*{\texorpdfstring{Appendices}{Appendices}}
\phantomsection\label{bp1}%
\addcontentsline{toc}{part}{\texorpdfstring{Appendices}{Appendices}}
% text: prose-PA-01
\phantomsection\label{bp2}%
To define \(Rf_!\) in Part I we must compactify every separated morphism of finite type, and to prove proper base change we must dominate every proper morphism by a projective one. These are Nagata's theorem and Chow's lemma (Appendix \bpout{A}). Nagata's theorem, followed by normalization, also gives the smooth compactification of a curve used in Part III (Lemma \bpout{14.1}).

Proper base change ends with a change of separably closed ground field. There we need the cohomology of curves for all torsion coefficients, including torsion divisible by the characteristic (Theorem \bpout{4.20}). Purity starts from sheaves of rank one. Their geometric monodromy is finite because Frobenius fixes no class in \(H^1\) of a curve (Proposition \bpout{14.3}). For coefficients prime to the characteristic, both facts come from the Jacobian. Indeed, \(H^1(C,\mu_n)\) is the \(n\)-torsion of the Jacobian, and a fixed class would give \(\ell\)-power torsion points of unbounded order among its finitely many rational points. Appendices \bpout{Appendix B}, \bpout{Appendix C} and \bpout{Appendix D} prove both facts.

Rankin's argument in Chapter \bpout{14} rests on the determinant weights of tensor and exterior powers (Proposition \bpout{14.19}). These come from the structure of the monodromy group (Lemma \bpout{14.15}), and Appendix \bpout{F} proves the facts about algebraic groups over \(\overline{\mathbf Q}_\ell\) behind it.

Appendix \bpout{E} is of a different kind. Following Varshavsky [\hyperref[bp3]{Var05}], it computes the local terms of contracting correspondences in a second way; Chapter \bpout{12} does not depend on it.


\clearpage\phantomsection\hypertarget{bp-outside}{}%
\section*{Outside this export}
\addcontentsline{toc}{section}{Outside this export}
References marked \ensuremath{{}^\ast} are to parts of the blueprint this document leaves out:
\begin{description}
\item[Theorem 4.20] Cohomology of curves, in Ⅰ The six operations (\bpslug{theorem-{}\allowbreak{}cohomology-{}\allowbreak{}curves})
\item[Chapter 12] The trace formula, in Ⅱ Duality and the trace formula (\bpslug{chap-{}\allowbreak{}The-{}\allowbreak{}trace-{}\allowbreak{}formula})
\item[Chapter 14] Purity of real sheaves, in Ⅲ Weights and the Riemann hypothesis (\bpslug{chap-{}\allowbreak{}Purity-{}\allowbreak{}of-{}\allowbreak{}real-{}\allowbreak{}sheaves})
\item[Lemma 14.1] A proper smooth model, in Ⅲ Weights and the Riemann hypothesis (\bpslug{lemma-{}\allowbreak{}proper-{}\allowbreak{}smooth-{}\allowbreak{}model})
\item[Proposition 14.3] No fixed class in degree one, in Ⅲ Weights and the Riemann hypothesis (\bpslug{proposition-{}\allowbreak{}no-{}\allowbreak{}fixed-{}\allowbreak{}class-{}\allowbreak{}degree-{}\allowbreak{}one})
\item[Lemma 14.15] The identity component, in Ⅲ Weights and the Riemann hypothesis (\bpslug{lemma-{}\allowbreak{}identity-{}\allowbreak{}component})
\item[Proposition 14.19] Determinant weights of tensor operations, in Ⅲ Weights and the Riemann hypothesis (\bpslug{proposition-{}\allowbreak{}determinant-{}\allowbreak{}weights-{}\allowbreak{}tensor-{}\allowbreak{}operations})
\item[Appendix A] Nagata compactification and Chow's lemma, in Appendices (\bpslug{chap-{}\allowbreak{}Nagata-{}\allowbreak{}compactification-{}\allowbreak{}and-{}\allowbreak{}Chows-{}\allowbreak{}lemma})
\item[Appendix B] Abelian varieties, in Appendices (\bpslug{chap-{}\allowbreak{}Abelian-{}\allowbreak{}varieties})
\item[Appendix C] The Picard scheme of a curve, in Appendices (\bpslug{chap-{}\allowbreak{}The-{}\allowbreak{}Picard-{}\allowbreak{}scheme-{}\allowbreak{}of-{}\allowbreak{}a-{}\allowbreak{}curve})
\item[Appendix D] Cohomology of torsion sheaves on curves, in Appendices (\bpslug{chap-{}\allowbreak{}Cohomology-{}\allowbreak{}of-{}\allowbreak{}torsion-{}\allowbreak{}sheaves-{}\allowbreak{}on-{}\allowbreak{}curves})
\item[Appendix E] Specialization and contracting correspondences, in Appendices (\bpslug{chap-{}\allowbreak{}Specialization-{}\allowbreak{}and-{}\allowbreak{}contracting-{}\allowbreak{}correspondences})
\item[Appendix F] Algebraic groups in characteristic zero, in Appendices (\bpslug{chap-{}\allowbreak{}Algebraic-{}\allowbreak{}groups-{}\allowbreak{}in-{}\allowbreak{}characteristic-{}\allowbreak{}zero})
\end{description}
\clearpage\phantomsection%
\section*{References}
\addcontentsline{toc}{section}{References}
\bpbibitem{BBD82}{bp4}{A. Beilinson, J. Bernstein, and P. Deligne. \emph{Faisceaux pervers}, volume 100 of \emph{Astérisque}. 1982.}

\bpbibitem{BS15}{bp5}{B. Bhatt and P. Scholze. The pro-étale topology for schemes. \emph{Astérisque}, 369:99–201, 2015. \href{https://people.mpim-bonn.mpg.de/scholze/proetale.pdf}{Author's manuscript}.}

\bpbibitem{Con07}{bp6}{B. Conrad. Deligne's notes on Nagata compactifications. \emph{Journal of the Ramanujan Mathematical Society}, 22:205–257, 2007. \href{https://math.stanford.edu/\~conrad/papers/nagatafinal.pdf}{Author's manuscript}.}

\bpbibitem{Del74}{bp7}{P. Deligne. La conjecture de Weil. I. \emph{Publ. Math. IHÉS}, 43:273–307, 1974.}

\bpbibitem{Del80}{bp8}{P. Deligne. La conjecture de Weil. II. \emph{Publ. Math. IHÉS}, 52:137–252, 1980.}

\bpbibitem{Den22}{bp9}{C. Deninger. There is no “Weil-”cohomology theory with real coefficients for arithmetic curves. \emph{Ann. Sc. Norm. Super. Pisa Cl. Sci.}, 2022. To appear. arXiv:\href{https://arxiv.org/abs/2204.02714}{2204.02714}.}

\bpbibitem{Deu41}{bp10}{M. Deuring. Die Typen der Multiplikatorenringe elliptischer Funktionenkörper. \emph{Abh. Math. Sem. Univ. Hamburg}, 14:197–272, 1941.}

\bpbibitem{Eke90}{bp11}{T. Ekedahl. On the adic formalism. In \emph{The Grothendieck Festschrift, Vol. II}, volume 87 of \emph{Progress in Mathematics}, pages 197–218. Birkhäuser, 1990.}

\bpbibitem{Has36}{bp12}{H. Hasse. Zur Theorie der abstrakten elliptischen Funktionenkörper. I, II, III. \emph{J. Reine Angew. Math.}, 175:55–62, 69–88, 193–208, 1936.}

\bpbibitem{Hum72}{bp13}{J. E. Humphreys. \emph{Introduction to Lie Algebras and Representation Theory}, volume 9 of \emph{Graduate Texts in Mathematics}. Springer, 1972.}

\bpbibitem{Ill94}{bp14}{L. Illusie. Autour du théorème de monodromie locale. \emph{Astérisque}, 223:9–57, 1994. §4. \href{https://www.numdam.org/item/AST_1994__223__9_0/}{Original article}.}

\bpbibitem{Kat88}{bp15}{N. M. Katz. Travaux de Laumon. \emph{Astérisque}, 161–162:105–132, 1988. Séminaire Bourbaki, exp. 691.}

\bpbibitem{KW01}{bp16}{R. Kiehl and R. Weissauer. \emph{Weil conjectures, perverse sheaves and \(\ell\)-adic Fourier transform}, volume 42 of \emph{Ergebnisse der Mathematik}. Springer, 2001.}

\bpbibitem{Lau87}{bp17}{G. Laumon. Transformation de Fourier, constantes d'équations fonctionnelles et conjecture de Weil. \emph{Publ. Math. IHÉS}, 65:131–210, 1987.}

\bpbibitem{Mil86a}{bp18}{J. S. Milne. \href{https://www.jmilne.org/math/xnotes/AVs.pdf}{Abelian Varieties}. In \emph{Arithmetic Geometry}. 1986. Corrected online version.}

\bpbibitem{Mil86b}{bp19}{J. S. Milne. \href{https://www.jmilne.org/math/xnotes/JVs.pdf}{Jacobian Varieties}. In \emph{Arithmetic Geometry}. 1986. Corrected online version.}

\bpbibitem{Ran39}{bp20}{R. A. Rankin. Contributions to the theory of Ramanujan's function \(\tau(n)\) and similar arithmetical functions. II. \emph{Proc. Cambridge Philos. Soc.}, 35:357–372, 1939.}

\bpbibitem{SGA4}{bp21}{M. Artin, A. Grothendieck, and J.-L. Verdier, editors. \emph{Théorie des topos et cohomologie étale des schémas}, volumes 269, 270, 305 of \emph{Lecture Notes in Mathematics}. Springer, 1972–1973. SGA 4.}

\bpbibitem{SGA4½}{bp22}{P. Deligne. \emph{Cohomologie étale}, volume 569 of \emph{Lecture Notes in Mathematics}. Springer, 1977. SGA 4½; \emph{Théorèmes de finitude}, §3, and \emph{Rapport sur la formule des traces}, §3, and \emph{Sommes trigonométriques}, §1.}

\bpbibitem{SGA5}{bp23}{A. Grothendieck, editor. \emph{Cohomologie \(\ell\)-adique et fonctions \(L\)}, volume 589 of \emph{Lecture Notes in Mathematics}. Springer, 1977. SGA 5.}

\bpbibitem{SGA7II}{bp24}{P. Deligne and N. Katz, editors. \emph{Groupes de monodromie en géométrie algébrique}, volume 340 of \emph{Lecture Notes in Mathematics}. Springer, 1973. SGA 7 II; Exposé XIII, (2.1.7).}

\bpbibitem{Stacks}{bp25}{The Stacks Project Authors. \href{https://stacks.math.columbia.edu}{The Stacks Project}. Individual results are cited by their stable tags in the text.}

\bpbibitem{Var05}{bp3}{Y. Varshavsky. Lefschetz–Verdier trace formula and a generalization of a theorem of Fujiwara. 2005. Version 2. arXiv:\href{https://arxiv.org/abs/math/0505564}{math/0505564}. \href{https://arxiv.org/abs/math/0505564v2}{Author's preprint}.}

\bpbibitem{Wei48a}{bp26}{A. Weil. \emph{Sur les courbes algébriques et les variétés qui s'en déduisent}. Hermann, Paris, 1948.}

\bpbibitem{Wei48b}{bp27}{A. Weil. \emph{Variétés abéliennes et courbes algébriques}. Hermann, Paris, 1948.}

\bpbibitem{Wei49}{bp28}{A. Weil. Numbers of solutions of equations in finite fields. \emph{Bull. Amer. Math. Soc.}, 55:497–508, 1949.}

\end{document}
